\section{Ethics for data science}


\subsection{Introduction}

\subsubsection{Law, moral and ethics}

Ethics is really old but the application in data science is fairly recent. Ethics is really important in fields where the consequences are really important, like healthcare. 

\begin{itemize}
    \item \textbf{Ethics} \ra Moral principles, branch of philosophy that studies the basis of moral. Set of rules of conduct. Aims to direct the individual towards good practices. Individual perfecting. Based on the relationship of the individual with himself. Ethical rules are filtered by the individual and it's conscience. Violation of ethical rules can lead to self-imposed sanctions. FOcussed on duties and obligations but not on rights.
    \item \textbf{Law} \ra Aims to regulate the essential relationships within a given society. Harmonious interactions and to solve conflicts between different people. Laws don't need individual conscience to be applied. Legal violations will lead to loss of liberties and life. Equal focus on duties and rights.
\end{itemize}

Laws waits for a conduct while ethics anticipate that conduct. Everyone can form an ethical opinion while legal ones are given by a judge. In ethical settings suspicion can affect someone while in legal ones people are innocent until proven guilty. 

\vspace{10pt}

Historically we saw innovations that lead to ethical principles and definitios (artificial life support created the legal definition of death or human genes discovery gave us the legal distinction of person vs thing).


\subsubsection{What is ethics}

We use ethics to make decisions related to choices that have moral dimensions. How one should act.

Ethics is made from two big branches

\begin{enumerate}
    \item Non-normative ethics \ra Studies what effectively the case is and not what ethically ought to be. We can have in this setting descriptive ethics (studies people beliefs about moralities) and meta-ethics (studies the meaning of moral propositions like right, obligations etc)
    \item Normative ethics \ra Studies the practical means of determining a moral course of action. It's applied ethics.
\end{enumerate}

We can also define the professional ethics branch, made of ethical gudielines set accepted by the majority of professionals. At a more granular level there exists also specific professional ethical rules. Some professions will adopt more detailed codes called Codes of ethical conduct. A public policy is a set of enforceable guidelines set by an official public body. 

\[\text{Public policy} = \text{Ethical principles} + \text{Empirical factors}\]



\subsubsection{Ethical dilemmas}

\begin{itemize}
    \item Psychologist and patient
    \item Cancer genes daughter
\end{itemize}


Ethical Dilemmas are circumstances in which ethical obligations demand or
appear to demand that a person adopt each of two (or more) alternative but
incompatible actions, such that the person cannot perform all the required
actions. Conflicts between ethical requirements and self-interest are
 practical dilemmas
and not necessarily
 ethical dilemmas

\vspace{10pt}

Tools of ethical decision making

\begin{itemize}
    \item \textbf{Basic ethics principles} \ra Aim to express general norms of common ethics (respect for autonomy, nonmaleficence, beneficence, justice)
    
    \item \textbf{Ethical rules} \ra More specific in content and more restricted in scope than principles. (however, both establish rights and obligations). Function as precise guides to action that direct us in each circumstance. They can be (substantive, authority and procedural)
    
    \item \textbf{Weighting and balancing} \ra Principles, rules, professional
obligations, and rights often
need to be weighted and
balanced.
Balancing is the process
of finding reasons to
support beliefs about
which ethical norms
should prevail. It is a process of practical astuteness, discriminating intelligence, and sympathetic
responsiveness. Necessary conditions are:
\begin{enumerate}
    \item  Good reasons can be offered to act on the overriding norm
rather than on the infringed
\item The moral objective justifying the infringement has a realistic
prospect of achievement

\item No morally preferable alternative actions are available
\item  The lowest level of infringement, commensurate with achieving
the primary goal of the action, has been selected
\item  Any negative effects of the infringement have been minimized
\item All affected parties have been treated impartially
\end{enumerate}

    \item \textbf{Virtues and emotions} \ra We want tho avoid ectremes situations like detachment or extreme attachment. In the absence of public and institutional constraints, partiality towards others is
ethically permissible and is the socially expected form of interaction. some necessary elements are discernment ( the ability to make fitting judgements and reach decisions without being influenced by extraneous factors), trustworthiness (Confident belief and reliance on the moral character and competence of another
person. Confidence that another will act with the right motives and in accordance with
ethical norms), integrity (Soundness, wholeness and integration of moral character. In a restricted sense means
fidelity in adherence to moral norms), concentiousness (Motivation to do what is right because it is right after trying with due diligence to
determine what is right) and compassion (Active regard for another´s welfare with an imaginative awareness and emotional
response of sympathy, tenderness and discomfort at another´s misfortune or suffering.)
\end{itemize}

\vspace{10pt}

Main ethical theories
\begin{itemize}
    \item Divine command
    \item Natural Law
    \item Virtue ethics
    \item Deontology
    \item Utilitarianism
\end{itemize}

A core point of ethical dilemmas is the disagreement that they can create given by different backgrounds, ideas or the quantity of information that everyone participating has. Disagreements may persist even among ethically committed persons.
Ethical disagreements may discourage persons who deal with practical problems but
should be no cause for skepticism about ethical thinking. 

\vspace{10pt}

Data ethics works on many levels based on the strenght of normative instruments. Remedies to violations can be
\begin{itemize}
    \item Prevention
    \item Detection/Investigation/Accusation
    \item Sanctions 
    \begin{itemize}
        \item Legal
        \item Disciplinary
        \item Professional
        \item Social
        \item Individual
    \end{itemize}
\end{itemize}

\subsection{History of data ethics}

Some case studies

\begin{itemize}
    \item Love is blind day
    \item Nurnmberg trial
    \item Tuskegee Experiment
    \item Milgram
    \item Rosenhan
    \item Tearoom trade
    \item Stanford prison 
\end{itemize}

Possible ways to approach digital revolution
\begin{itemize}
    \item Unchecked optimism
    \item Go backwards
    \item Accept with uneasiness
    \item Bend it towards justice and fairness
\end{itemize}

An important element for the last point is informed consent created after the Nuremberg code, it says that voluntary consent of the human subject is absolutely essential. Also the degree of risk should be proportional to the humanitarian impact.
Informed consent is made of 

\begin{itemize}
    \item Previous and precise aknowledgement of the situation
    \item Freedom from coercion
    \item Competence
    \item Flexibility
    \item Attention to vulnerable groups
\end{itemize}

Sometimes consent is impossible to obtain but it must be regularly evaluated.
Consent must be
\begin{itemize}
    \item Informed
    \item Voluntary
    \item Ethical 
\end{itemize}

The data ethics framewok is made of 

\begin{itemize}
    \item Overarching principles \ra Transparency, accountability, fairness
    \item Specific actions
\end{itemize}  


\subsection{Basics of organizational ethics}

Some snippets from Antonio Martins life. We learn our names and origins. 

\vspace{10pt}

Ethics is a set of principles, values and standards of conduct which guides individuals and organizations determining whats right and wrong. It involves making decisions and taking actions that are morally right considering the impact on stakeholders. Video on ethics.

\vspace{10pt}

We can consider organizational ethics as 
\begin{itemize}
    \item Trust and reputation
    \item Legal compliance
    \item Employee morale and engagement
\end{itemize}

We live in a society (professor said it). Joao from Leiria talks. 

\vspace{10pt}

Why companies need to be ethical

\begin{itemize}
    \item Building trust and reputation 
    \item Enhancing employee engagement and retention
    \item Mitigating risks and legal liabilities
\end{itemize}

Area where to use ethical approaches

\begin{itemize}
    \item Improving decision making and innovation
    \item Maintaining stakeholder relationships
    \item Demonstrating corporate social responsibility
\end{itemize}

When we talk about ethichs we talk about frameworks for organizational behaviour thus shaping how relationships work in an organization. Prioritizing ethical principles will make organizations creating sustainable and socially responsible business environment. When can also talk about ethical conduct, so how the company actually acts, but is not strictly morally right. 

\vspace{10pt}

Some ethical theories

\begin{enumerate}
    \item \textbf{Utilitarianism} \ra Maximise the happines of as many people as possible
    \item \textbf{Deontology} \ra We have some moral duties to follow
    \item \textbf{Virtue ethics} \ra Focus on virtues will lead to morally good lives
\end{enumerate}

Excercise about questions. 

\vspace{10pt}

Ethical stakeholders
\begin{itemize}
    \item Individuals
    \item Organizations
    \item Society
\end{itemize}

Key issues in data ethics 

\begin{itemize}
    \item Privacy
    \item Fairness
    \item Accountability
    \item Emerging challenges
\end{itemize}

\subsubsection{Privacy and confidentiality}

Personal data is a cornestone in digital economies, it's misuse can lead to privacy violations and heavy harm on all the actors invoved. Two important ways to save these data are a\textbf{Anonymization} and \textbf{differential privacy}. We define personal data as any information that relates to an identified or identifiable individual. Like basic identifiers, sensitive data, behavioural data and location data. 

Personal data have many values
\begin{itemize}
    \item Economic
    \item Strategic
    \item Technological advancement
    \item Societal insights
    \item Personal 
\end{itemize}

The main beneficiaries are
\begin{itemize}
    \item Businesses
    \item Governments
    \item Individuals
\end{itemize}

The risks are 
\begin{itemize}
    \item Data breaches
    \item Identity theft
    \item Surveillance and data erosion
    \item Exploitation
    \item Lack of consent
\end{itemize}

Two example of frameworks to save data privacy can be the \textbf{GDPR} and the \textbf{CCPA}


\subsubsection{Bias and fairness in statistical models}

\subsubsection{Data sharing and collaboration}

\subsubsection{Responsible use of predictive analytics and machine learning}

\subsubsection{Transparency and accountability}

\subsubsection{Ethical decision making process}




























